\section*{Devoir maison }

On s'intéresse au code suivant :

\begin{lstlisting}[language=Python]
class Recette:

    def __init__(self, nom, ingredients, quantites, unites):
        self.ingredients = ingredients
        self.quantites = quantites
        self.unites = unites
        self.nom = nom
    
    def ajouter(............................):
        ....................................................
        ....................................................
        ....................................................

    def calcul_prix(self, prix_unitaires):
        ''' Calcule le prix de la recette a partir des prix unitaires.  '''
        total = 0
        for i in range(len(self.ingredients)):
            total += prix_unitaires[i] * self.quantites[i]
        return total

    def liste_course(self, cellier):
        ''' Retourne une liste de tuple (ingredient, quantite).'''
        pass
                       
\end{lstlisting}

\begin{Exercise}[title={POO - Recette}]
    \Question Lister les attributs de la classe Recette en précisant leur type.\\
    \emptybox{1cm}
    \Question Compléter la méthode ajouter de la classe Recette qui permet d’ajouter un ingrédient, sa quantité et son unité aux listes correspondantes.
    \Question Quel est le type du renvoi des méthodes ajouter, calcul\_prix et liste\_course ?\\
    \emptybox{1cm}
    \Question Écrire l’instruction permettant d’instancier l’objet tartiflette caractérisant la recette de la tartiflette avec 	500 g de reblochon, 200 g d'oignon émincés, 200 g de lardons fumés, 1 kg de pomme de terre.\\
    \emptybox{1cm}
    \Question On se raconte qu'il vaut mieux mettre 1kg de reblochon dans la recette de la tartiflette écrire les instructions permettant de mettre à jour la recette tartiflette.\\
    \emptybox{2cm}
\begin{lstlisting}[language=Python]
muffin = Recette("Muffin au chocolat",
                 ["farine", "chocolat", "lait", "oeuf"],
                 [100, 50, 200, 2],
                 ["grammes", "grammes", "grammes", "unites"])
prix = muffin.calcul_prix([0.1, 1.1, 0.75, 0.25])
print(f"Le prix de vente du {muffin.nom} est de {prix} €.")
\end{lstlisting}
    \Question Sachant que le prix de muffin.calcul\_prix([0.1, 1.1, 0.75, 0.25]) est de 215.5.  Qu'affiche exactement la ligne 6 du code ci-dessus ?\\
    \emptybox{1cm}
    \Question Donner le tableau d'avancement tour par tour de l'appel de la méthode calcul\_prix de la ligne 5 du code ci-dessous.\\
    \emptybox{5cm}
    \Question Recopier et compléter la méthode liste\_course qui retourne la liste de course qui est une liste de tuple (ingrédient, quantité).\\
    Un ingrédient est dans la liste de course si :
    \begin{itemize}
        \item Il n'apparait pas dans le cellier, dans ce cas, on prend la quantité de la recette.
        \item La quantité dans le cellier est inférieur à celle de la recette, dans ce cas, on prend la différence entre la quantité de la recette et celle du cellier.
    \end{itemize}
    \emptybox{10cm}
\end{Exercise}